\documentclass[11pt,a4paper]{article}
\usepackage[margin=25mm]{geometry}
\usepackage{amsmath,amssymb,amsthm,hyperref}
\setlength{\emergencystretch}{2em}
\newtheorem{lemma}{Lemma}
\newtheorem{theorem}[lemma]{Theorem}
\newtheorem{corollary}[lemma]{Corollary}
\newcommand{\gen}[1]{\langle#1\rangle}
\title{An orbit deficit for ambient-eligible tail cuts}
\author{Human Compression project: L044 working continuation}
\date{}
\begin{document}\maketitle
Let $G=\gen{\rho_0,\ldots,\rho_{r-1}}$ be a transitive
string permutation group generated by involutions. Write
$G_j=\gen{\rho_k:k\ne j}$, and call a $j$-edge
\emph{ambient-eligible} when its endpoints lie in different
$G_j$-orbits. Fix $0\leq t<r-1$, put
\[
 I=\{t+1,\ldots,r-1\},\qquad H=G_{>t},\qquad
 U=\operatorname{supp}(H),\qquad d=|I|.
\]
Suppose there are at least two ambient-eligible edges of
each colour in $I$. A tail selection $Q$ chooses exactly
two such edges per colour, on vertex set $U$.
All actual edges of these colours lie in $U$ and inside
individual $H$-orbits. Full IP, evenness and connectedness
of the noncommuting diagram are not assumed.

The paid paired-cut square-normal-form proof applies to this
interval of colours: its orbit-cut exchanges use ambient
eligibility, and every global middle-colour move has its
middle index in $I$. Choose a square-only selection with
minimum number $q$ of squares. As in that proof, \emph{every}
such minimum selection has at most one square per component:
otherwise minimize the distance between two squares over all
minimum selections and apply the incident, ordinary or global
middle move. Thus every component is a tree or is unicyclic
with a square. No transitivity of $H$ is used here.

\begin{lemma}[Localisation and ordinary slides]\label{local}
If $S$ is the square in a component $C$, every
ambient-eligible edge of a square colour has both endpoints
in $C$. An ordinary square slide along a selected bridge
leaving $S$ preserves the vertex set of $C$ and all other
component vertex sets.
\end{lemma}
\begin{proof}
Let $a$ be a square colour. An eligible $a$-edge cannot
have both endpoints in another component: its selected
connecting path would avoid $a$. If its endpoints belonged
to different components, delete an $a$-edge of $S$ and
insert it. The deletion destroys $S$ without disconnecting
$C$, and the insertion joins components without creating
a cycle. This contradicts minimum $q$ and proves localisation.

For the slide write its first bridge as
$v\mathrel{-_\ell}w$, with $v\in S$ and $\ell$ commuting
with a square colour $a$. Put $v'=v^{\rho_a}$ and
$w'=w^{\rho_a}=(v')^{\rho_\ell}$. The completing edges
$ww'$ and $v'w'$ are eligible by the orbit-cut exchange
lemma. Localisation puts $w,w'$ in $C$.
Delete the opposite $a$-edge of $S$, insert $ww'$, and
complete $v'w'$ if missing by deleting the other selected
$\ell$-edge. All added edges lie in the new full square,
so the square-block lemma permits no other newly created
cycle. The old square is destroyed, and no other old cycle
can contain an $\ell$-edge: the first selected $\ell$-edge
is a bridge, while a cycle would require both selected edges
of that colour.

If the parallel $\ell$-edge was already selected, deleting
the old square edge leaves a connected tree on $C$ and
inserting the new edge restores one cycle. Otherwise the
deleted other $\ell$-edge cannot lie outside $C$: after
deleting the square edge, the two added internal edges would
give the still connected $C$ two cycles, although its only
possible cycle is the new square. When the deleted edge is
in $C$, its edge count remains $|C|$, and its only cycle is
the new square; hence it is connected. All other components
are untouched, proving the partition assertion.
\end{proof}

Fix an $H$-orbit $O\subseteq U$ and suppose some minimum
selection has no tree component in $O$. Equivalently its
number $E_O$ of selected edges in $O$ equals $|O|$.
The following statements apply to every minimum selection
having this equality.

\begin{lemma}[Neighbours and middle transport inside a full orbit]
\label{middle}
Every full-graph neighbour of a vertex of a square $S$
in a component $C\subseteq O$ lies in $C$. A global
middle-colour move at $S$ preserves $E_O=|O|$ and transports
the component partition within $O$ by its middle generator.
\end{lemma}
\begin{proof}
Write the square colours as $a,b$. A full neighbour of colour
$a$ or $b$ lies in $S$. If its colour $k$ commutes with $a$,
conjugating the selected $a$-edge by $\rho_k\in G_a$
gives an actual ambient-eligible $a$-edge. Lemma~\ref{local}
puts both endpoints, including that neighbour, in $C$.
The same argument applies when $k$ commutes with $b$.

The only remaining case is $a=i-1,b=i+1,k=i$. Both square
colours lie in $I$, so $i\in I$ and $\rho_i\in H$.
Retain the four edges of $S$ and transport every other
selected edge by $\rho_i$, the accepted global middle move.
This remains a minimum square-only selection. All $H$-orbits
are preserved by $\rho_i$, and all four retained edges lie
in $O$, so $E_O=|O|$ is preserved. Since every component of
a minimum selection has at most one cycle, that equality
forces every component within $O$ to be unicyclic.

As an operation on $Q^{\rho_i}$, the move deletes
$S^{\rho_i}$ and inserts $S$. The deletion splits
$C^{\rho_i}$ into four trees. Every other component in $O$
remains unicyclic. An inserted square vertex cannot lie in
one of those other components, since it would then join
two squares in one component. Hence all four inserted vertices
lie in $C^{\rho_i}$. Inserting $S$ reconnects all four
pieces, since otherwise a tree in $O$ would remain.
Thus the partition within $O$ is exactly the transported
partition. Also $S\subseteq C^{\rho_i}$, equivalently
$S^{\rho_i}\subseteq C$, proving the remaining neighbour
assertion. Components with trees in other $H$-orbits cannot
interfere: the middle generator preserves each orbit.
\end{proof}

\begin{theorem}[A tree in every proper tail orbit]\label{tree}
Every proper $H$-orbit $O\subsetneq\Omega$ in $U$ has a
tree component in a minimum-square tail selection.
Consequently, if $H$ is intransitive on $\Omega$ and has
$c$ nontrivial orbits, then
\[
                         2d\leq |U|-c.                 \tag{1}
\]
\end{theorem}
\begin{proof}
Suppose $O$ has no tree. Consider all minimum-square
selections with $E_O=|O|$. By transitivity of the full $G$
and properness of $O$, a full edge leaves $O$. Its inside
endpoint lies in a selected component $C\subseteq O$, which
has a unique square $S$. A selected path from $S$ to that
endpoint followed by the exit edge gives a connector.
Minimize its length over all such selections, squares,
selected paths and exit edges. The exit colour $x$ satisfies
$x\leq t$, since every generator in $I$ preserves $O$.
By Lemma~\ref{middle} the selected path is nonempty.

Let its first edge be $v\mathrel{-_\ell}w$, $v\in S$.
If $\ell$ commutes with a square colour $a$, use the
ordinary slide. Its other selected $\ell$-edge cannot be
farther along the connector's selected path. Indeed delete
the first bridge and let $W$ be the side containing $w$.
Both selected $a$-edges lie outside $W$. If
$w'=w^{\rho_a}$ were in $W$, its path from $w$ would
avoid $a$, contradicting eligibility. If it were in
$C\setminus W$, its path from $v^{\rho_a}$ would avoid
$\ell$, since both selected $\ell$-edges are on or beyond
the deleted bridge. This too contradicts eligibility.
Thus this repeated-colour possibility would put $w'$ outside
$C$, contrary to Lemma~\ref{local}.
The slide therefore preserves the rest of the selected path,
the exit edge (whose colour is outside $I$), and the orbit
edge count. Its new square contains $w$, giving a shorter
connector, a contradiction.

Otherwise $S$ has colours $i-1,i+1$ and $\ell=i$.
Here $i\geq t+2$, so $\rho_i$ commutes with the exit
generator $\rho_x$, since $x\leq t$.
Perform the global middle move. Every selected edge of the
connector after the first is outside $S$ and therefore
transports to a selected edge. The full exit edge transports
to an actual $x$-edge. Its outside endpoint remains outside
$O$, because $\rho_i\in H$ preserves $O$ and its complement.
The transported path begins at $w^{\rho_i}=v\in S$.
Together with the transported exit it gives a connector
shorter by one, still in a minimum selection with
$E_O=|O|$. Possible vertex coincidences only shorten it
further. This is the final contradiction.

Thus each proper nontrivial $H$-orbit contains a tree.
Every tree component contributes one to vertices minus edges,
and each unicyclic component contributes zero. Summing these
deficits over the $c$ nontrivial orbits gives (1).
\end{proof}

\begin{corollary}[The intransitive-tail rank step]
Suppose $H$ has at least two nontrivial orbits on $U$ and
$t\leq(n-|U|-1)/2$. Under the ambient-eligibility hypotheses
above, $r\leq\lfloor(n-1)/2\rfloor$.
\end{corollary}
\begin{proof}
Equation (1) gives $2(r-1-t)\leq |U|-2$, hence
$2r\leq |U|+2t\leq n-1$.
\end{proof}

\paragraph{Source, attribution and remaining boundary.}
This is a proposed replacement for the \emph{ambient-eligible}
use of Cameron--Fernandes--Leemans--Mixer, \emph{Highest rank
of a polytope for $A_n$}, arXiv:1605.09173v1,
Proposition~5.27, physical pages~34--36. Its literal statement
assumes a paired fracture graph for $H$ itself; that means
crossing orbits of $H_j$, a weaker condition than crossing
ambient $G_j$-orbits. The proof above does not identify these
conditions. Later applications must supply two ambient-eligible
edges of each tail colour, or take their separate unique-crossing
branch. The broader literal local-eligibility assertion remains
separate, not refuted by this distinction alone.

The fracture framework and rank target are the source's.
The whole accepted paired-cut proof and equality-localisation
proof remain paid dependencies. The global middle-colour move
was supplied by the root collaborator and distinctly reviewed;
the orbit-deficit argument here is a new proof adaptation of
that accepted mechanism. No finite census, formal theorem,
whole L044 recovery, original or question credit is asserted.
\end{document}
